\vspace{0.5em}\noindent\textbf{Self-evaluation}\par

\vspace{0.5em}\noindent\textbf{Assessment summary}\par

This is my draft self-assessment based on the Internship Application
Tracker project. It is not an official supervisor rating.

During the project, I worked on a cloud-native internship application
platform that combines a React/Vite frontend, FastAPI backend, Node.js
Socket.IO chat service, PostgreSQL, DynamoDB, Redis, SQS, Lambda, S3,
CloudFront, EKS, and SageMaker. The most important learning outcome for
me was understanding how application code, Kubernetes, managed AWS
services, IAM, networking, CI/CD, and runtime evidence connect in a real
deployment.

\vspace{0.5em}\noindent\textbf{Skill rating table}\par

{
\begingroup
\small
\setlength{\tabcolsep}{3pt}
\renewcommand{\arraystretch}{1.15}
\begin{longtable}{@{}
  >{\raggedright\arraybackslash}p{(\linewidth - 4\tabcolsep) * \real{0.3000}}
  >{\raggedleft\arraybackslash}p{(\linewidth - 4\tabcolsep) * \real{0.4000}}
  >{\raggedright\arraybackslash}p{(\linewidth - 4\tabcolsep) * \real{0.3000}}@{}}
\toprule\noalign{}
\begin{minipage}[b]{\linewidth}\raggedright
Skill area
\end{minipage} & \begin{minipage}[b]{\linewidth}\raggedleft
Draft rating
\end{minipage} & \begin{minipage}[b]{\linewidth}\raggedright
Evidence and reflection
\end{minipage} \\
\midrule\noalign{}
\endhead
\bottomrule\noalign{}
\endlastfoot
AWS architecture & 8/10 & I helped shape the final architecture where
CloudFront routes static frontend traffic to S3 and dynamic
\texttt{/\allowbreak{}api}, \texttt{/\allowbreak{}chat}, and \texttt{/\allowbreak{}socket.\allowbreak{}io} traffic to the
ALB. I learned to distinguish proposal diagrams from deployed
evidence. \\
Kubernetes & 8/10 & I worked with EKS workloads for backend, chat,
outbox dispatcher, processing worker, and ai-service, including
Deployments, Services, readiness/liveness probes, HPA, and PDB. \\
Networking & 7/10 & I investigated issues such as a NAT Gateway
blackhole that prevented EKS nodes from joining and the AWS Load
Balancer Controller VPC ID failure. I still need more practice designing
VPCs from scratch. \\
IAM and security & 7/10 & I used GitHub OIDC, IRSA, runtime roles,
secret separation, S3 private access, SQS SSE, and idempotency. I need
to improve formal IAM policy review and permission boundary design. \\
CI/CD & 8/10 & I worked with GitHub Actions workflow modes such as
\texttt{validate}, \texttt{deploy-\allowbreak{}app}, \texttt{deploy-\allowbreak{}public},
\texttt{deploy-\allowbreak{}frontend}, \texttt{rollout}, and \texttt{full}, including
ECR SHA image verification and frontend S3 deployment. \\
Database and messaging & 8/10 & I learned how PostgreSQL stores business
data, processing jobs, idempotency records, and transactional outbox
events, while SQS carries committed events to Lambda and DynamoDB
supports chat and dedupe. \\
Serverless & 7/10 & I integrated the SQS-to-Lambda notification flow
conceptually and documented the successful Lambda smoke test, DynamoDB
dedupe, S3 archive, SES result, and partial batch failure behavior. \\
AI integration & 7/10 & I worked with the processing worker and
ai-service adapter pattern, where the worker keeps stable routes and the
adapter calls SageMaker endpoint \texttt{internship-\allowbreak{}qwen3-\allowbreak{}4b}. I need
more experience operating model endpoints cost-effectively. \\
Troubleshooting & 8/10 & I practiced evidence-led debugging for NAT, ALB
controller, GitHub Actions job conditions, SQS queue naming, CloudFront
routing, and DynamoDB reserved keyword errors. \\
Documentation & 8/10 & I converted implementation details and runtime
context into workshop, proposal, cost, security, testing,
troubleshooting, and cleanup documentation. \\
\end{longtable}
\endgroup
}

\vspace{0.5em}\noindent\textbf{Strengths}\par

\begin{itemize}
\tightlist
\item
  I can connect source code, deployment manifests, workflow files, and
  runtime evidence into one coherent architecture explanation.
\item
  I became more careful about distinguishing implemented code from
  verified production behavior.
\item
  I improved at troubleshooting AWS integration failures by starting
  from logs and command output instead of assumptions.
\item
  I learned how to document deployment procedures with warnings,
  expected results, and common errors.
\item
  I understand why frontend static hosting, EKS services, database
  transactions, SQS delivery, Lambda idempotency, and SageMaker
  inference belong in different parts of the architecture.
\end{itemize}

\vspace{0.5em}\noindent\textbf{Challenges}\par

\begin{itemize}
\tightlist
\item
  AWS service interactions can fail for reasons outside application
  code, especially IAM, route tables, controller configuration, and
  GitHub Actions variable scope.
\item
  Debugging EKS required understanding both Kubernetes objects and
  AWS-created resources such as ALB target groups and ENIs.
\item
  Cost estimation was difficult because exact AWS prices depend on live
  instance types, endpoint uptime, data transfer, and log retention.
\item
  AI integration required preserving the worker-facing contract while
  moving implementation details into the SageMaker adapter.
\item
  Documentation required discipline because it is easy to overstate
  something that exists in code but has not been verified in production.
\end{itemize}

\vspace{0.5em}\noindent\textbf{Lessons learned}\par

\begin{itemize}
\tightlist
\item
  Runtime evidence should be treated as stronger than old manifests when
  describing current production architecture.
\item
  A Kubernetes Service load-balances traffic but does not serialize
  database writes; correctness must come from constraints, idempotency,
  conditional writes, and transactions.
\item
  SQS Standard provides at-least-once delivery, so consumers must
  deduplicate events.
\item
  Lambda is useful for short event-driven processing, but it only
  reduces cost when it removes or scales down always-on capacity.
\item
  Static frontend delivery through S3 and CloudFront is a better fit
  than running a frontend pod in EKS for this project.
\item
  Worker workloads should be enabled only after their dependencies,
  especially SageMaker, are ready.
\end{itemize}

\vspace{0.5em}\noindent\textbf{Areas for improvement}\par

\begin{itemize}
\tightlist
\item
  Practice writing IAM policies with least privilege from the beginning
  instead of tightening them after deployment failures.
\item
  Learn more about VPC endpoint strategy to reduce NAT Gateway cost and
  dependency.
\item
  Add stronger production evidence collection for alarms, backups,
  encryption settings, and PITR.
\item
  Build a repeatable frontend/browser E2E test harness for login, job
  apply, chat, and AI processing flows.
\item
  Improve SageMaker cost control by testing scheduled shutdown,
  asynchronous inference, or lower-cost inference options.
\end{itemize}

\vspace{0.5em}\noindent\textbf{Future plan}\par

\begin{enumerate}
\def\labelenumi{\arabic{enumi}.}
\tightlist
\item
  Strengthen AWS networking knowledge, especially VPC routing, NAT,
  endpoints, security groups, and private service access.
\item
  Practice production IAM design with OIDC, IRSA, Lambda roles, and
  SageMaker execution roles.
\item
  Build a complete evidence package for deployed systems: CLI exports,
  logs, screenshots, Cost Explorer, and Pricing Calculator results.
\item
  Add automated E2E tests and smoke tests for the main candidate and HR
  flows.
\item
  Continue improving event-driven design with SQS, Lambda, DLQ handling,
  idempotency, and observability.
\end{enumerate}

\vspace{0.5em}\noindent\textbf{Career orientation}\par

This project confirmed that I want to keep developing skills in cloud
engineering, backend systems, DevOps, and AI-enabled applications. I am
especially interested in work that combines application development with
reliable infrastructure, security, automation, and practical
troubleshooting.
